% Template created by Jon Hood

\documentclass[letterpaper, 10pt, twoside]{article}
\usepackage{noto}
\usepackage{fancyhdr}
\usepackage{graphicx}
\usepackage{multirow}
\usepackage[table]{xcolor}
\usepackage[breaklinks=true]{hyperref}
\usepackage[letterpaper, top=1in, bottom=1in, left=1.5in, right=1in, includeheadfoot]{geometry}
\usepackage{wrapfig}

%                Main Variables
\newcommand{\repdate}{\formatdate{15}{5}{2020}}
\def \ProjectName{Security Technical Implementation Guide Qt Viewer}
\def \ProjectAcronym{STIGQter}
\def \ProjectVersion{1.0.1}

%-------------------------------------------------------------------------------
%                Header & Footer Setup
%-------------------------------------------------------------------------------
\pagestyle{fancy}
\lhead[\thepage]{\includegraphics[height = 2em]{images/STIGQter.pdf}}
\chead{\ProjectName\ --- \ProjectVersion}
\rhead[{\includegraphics[height = 2em]{images/STIGQter.pdf}}]{\thepage}
\cfoot{\ProjectAcronym\ \ProjectVersion}

\fancypagestyle{blank}
{
	\lhead{}
	\chead{}
	\rhead{}
	\cfoot{}
	\lfoot{}
}

%-------------------------------------------------------------------------------
%                Revisions Table
%-------------------------------------------------------------------------------
\definecolor{STIGQterBlue}{RGB}{30,72,124}
\newcounter{RevisionCounter}
\newenvironment{Revision}{
	\begin{center}
		\begin{tabular}{ | c | c | c | p{23em} | }
			\hline
			\multicolumn{4}{| c |}{\cellcolor{STIGQterBlue}\textbf{\textcolor{white}{Revision History}}} \\
			\hline
			\rowcolor{lightgray}
			\textbf{Date} & \textbf{Revision} & \textbf{Revised By} & \textbf{Reason} \\
			\hline
		}{
	\end{tabular}
\end{center}
}
\newcommand{\RevisionEntry}[3]{\stepcounter{RevisionCounter}
	#1 & \Alph{RevisionCounter} & #2 & #3 \\
	\hline}

%-------------------------------------------------------------------------------
%                Title Page
%-------------------------------------------------------------------------------
\newcommand{\headerlogo}{
	\includegraphics[width=.3\linewidth]{images/STIGQter.pdf}\\
	\vspace{.5em}
}
\newcommand{\docline}{\textmd{\textbf{STIGQter System and Services Acquisition (SA) Documentation\\}}}
\newcommand{\generator}{Generated By: Jon Hood}
\title{
	\headerlogo
	\docline
	\vspace{.5em}
	\normalsize{\generator}
}

\author{Jon Hood}
\date{\repdate}

%-------------------------------------------------------------------------------
%                PDF metadata
%-------------------------------------------------------------------------------
\hypersetup
{
	pdfauthor= (Jon Hood),
	pdftitle =  (STIGQter System and Services Acquisition)
}

\usepackage{attachfile}
\usepackage[backend=biber]{biblatex}
\usepackage{import}
\usepackage{graphicx}
\usepackage[utf8]{inputenc}
\usepackage{url}
\addbibresource{sources.bib}
\usepackage{setspace}
\usepackage{array}
\usepackage{booktabs}
\newcolumntype{L}{@{}>{\kern\tabcolsep}l<{\kern\tabcolsep}}
\usepackage{colortbl}
\usepackage{xcolor}
\usepackage{textcomp}
\usepackage{outlines}
\usepackage{setspace}
\usepackage{longtable}
\usepackage{enumitem}
\usepackage{listings}
\usepackage{pgfplots}
\usepgfplotslibrary{fillbetween}
\usepackage{datetime}
\newdateformat{changelog}{\THEMONTH/\THEDAY/\THEYEAR}

\lstdefinestyle{CStyle} {language=C}
\lstdefinestyle{CSharpStyle} {language=[Sharp]C}
\lstdefinestyle{PHPStyle} {language=php}
\lstdefinestyle{JavaStyle} {language=java}
\lstdefinestyle{BashStyle} {language=bash}
\lstdefinestyle{SqlStyle} {language=SQL}

\lstset{language=C}
\lstset{language=[Sharp]C}
\lstset{language=php}
\lstset{language=java}
\lstset{language=bash}
\lstset{language=SQL}

\newcolumntype{L}[1]{>{\raggedright\let\newline\\\arraybackslash\hspace{0pt}}m{#1}}
\newcolumntype{C}[1]{>{\centering\let\newline\\\arraybackslash\hspace{0pt}}m{#1}}
\newcolumntype{R}[1]{>{\raggedleft\let\newline\\\arraybackslash\hspace{0pt}}m{#1}}

%start the document
\begin{document}

%generate a title page
\thispagestyle{blank}
\maketitle
\newpage

\thispagestyle{blank}
This page intentionally left blank for 2-sided printing compatibility.
\newpage

%tell the page counter to restart
\setcounter{page}{1}
\pagenumbering{roman} %use lowercase Roman numerals for page numbers

%Revision History page
\begin{Revision}
\RevisionEntry{\date{\changelog\formatdate{16}{1}{2020}}}{Jon Hood}{Initial Release of SA Policy}
\RevisionEntry{\date{\changelog\formatdate{28}{1}{2020}}}{Jon Hood}{Prepare for 1.0 Release}
\RevisionEntry{\date{\changelog\formatdate{26}{2}{2020}}}{James Shelton}{Update to support not-applicable CCI documentation}
\end{Revision}
\newpage

%generate a table of contents
\setcounter{tocdepth}{2}
\tableofcontents
\cleardoublepage % go to next right-side page

%reset the page counter for the regular pages
\setcounter{page}{1}
\doublespacing
\pagenumbering{arabic} %use standard numbers for the page number

\section{Background}

STIGQter is an open source tool used to perform RMF and STIG compliance mapping. As an open-source application, SA is performed by the STIGQter administrators on GitHub at \url{https://github.com/squinky86/STIGQter}.

\section{Purpose}

This documentation is meant to satisfy System and Services Acquisition (SA) requirements for RMF and detail what pieces of a SA policy may be covered by this application. SA requirements of roles (Section~\ref{sec:roles}) and their responsibilities (Section~\ref{sec:responsibilities}) provide integrity to the application.

\section{Scope}

The SA policy describes development process work and procedures.

\section{Roles}
\label{sec:roles}

The following role(s) are defined as being applicable to STIGQter execution:
\begin{enumerate}
	\item \textbf{User} -- The user of the software is defined as the individual account given permissions to execute the software.
\end{enumerate}

The following role(s) are defined as being applicable to STIGQter GitHub administration:
\begin{enumerate}
	\item \textbf{Administrator} -- An administrator of the STIGQter GitHub project
\end{enumerate}

\section{Responsibilities}
\label{sec:responsibilities}

Responsibilities are described in the GPL version 3.

\section{Management Commitment}

STIGQter development is committed to occur in an approved open source repository for maximum visibility and distribution.

\section{Coordination}

No coordination with other organizations is required to operate this software.

\section{RMF Control Compliance}

\subsection{SA-4 -- Acquisition Process}

\subsubsection{SA-4(5) -- System / Component / Service Configurations}

\paragraph{Applicable CCIs:} CCI-3109, CCI-3110, CCI-3111

The application is delivered in a secure state and with no additional user accounts created.

\subsection{SA-5 -- Information System Documentation}

\paragraph{Applicable CCIs:} CCI-642, CCI-3124, CCI-3125, CCI-3126, CCI-3127, CCI-3128, CCI-3129, CCI-3130, CCI-3131, CCI-3132, CCI-3133, CCI-3134, CCI-3135, CCI-3136

Application Configuration Guide is included as part of the Configuration, Installation, and Usage Guide. All documentation is maintained in the project's GitHub repository.

\subsection{SA-10 -- Developer Configuration Management}

\subsubsection{SA-10(1) -- Software / Firmware Integrity Verification}

\paragraph{Applicable CCIs:} CCI-698

The software is hashed with SHA512 and Blake2b512. Hashes are included with all final releases.

\subsection{SA-11 -- Developer Security Testing and Evaluation}
\label{sec:sa-11}

\paragraph{Applicable CCIs:} CCI-3171, CCI-3172, CCI-3173, CCI-3174, CCI-3175, CCI-3176, CCI-3177, CCI-3178

Scans are performed, at a minimum, on each release of the software. Most scans are performed with each commit of the source code using several automated tools. For each commit, Travis-CI is initialized and begins a scanning, debug build of the software. The following tools are run automatically:
\begin{itemize}
	\item Codacity - Code Quality Scans
	\item Codecov - Code Coverage Analysis
	\item CodeFactor - Code Quality Scans
	\item Coveralls - Code Coverage Analysis
	\item GuardRails - Code Security Scans
	\item Kritika - Code Security and Coverage
	\item lgtm - Code Security and Quality Scans
	\item SonarCloud - Code Security and Quality Scans
	\item SourceLevel - Code Security Scans
\end{itemize}

Additionally, each release of the application is scanned by Coverity, Clazy, and Clang-Analyzer for defects.

When quality or security defects are identified, they are resolved before the next release of the application. An application will not be released if high or critical risk findings are detected in the security or quality scans.

\subsubsection{SA-11(2) -- Threat and Vulnerability Analyses}

\paragraph{Applicable CCIs:} CCI-3182

Testing is the responsibility of the Administrators for the project on GitHub. Tests are run by Travis-CI for every new commit.

\subsubsection{SA-11(3) -- Independent Verification of Assessment Plans / Evidence}

\paragraph{Applicable CCIs:} CCI-3183, CCI-3184, CCI-3185, CCI-3186

The scanning services listed in Section~\ref{sec:sa-11} are independent of the development team, publish their assessment results publicly, and have been granted full access to the source repositories.

\subsubsection{SA-11(4) -- Manual Code Reviews}

\paragraph{Applicable CCIs:} CCI-3187, CCI-3188

The GitHub project administrators are responsible for verifying that the scanning documented in Section~\ref{sec:sa-11} is performed. Code coverage of security scans is 100\%.

\subsubsection{SA-11(8) -- Dynamic Code Analysis}

\paragraph{Applicable CCIs:} CCI-3197

High and Critical risk security findings from the scans in Section~\ref{sec:sa-11} are tracked in the project's GitHub issue tracker. Some historical data is preserved in the cloud scanner's history.

\subsection{SA-15 -- Development Process, Standards, and Tools}

\paragraph{Applicable CCIs:} CCI-3233, CCI-3234, CCI-3235, CCI-3236, CCI-3237, CCI-3238, CCI-3239, CCI-3240, CCI-3241, CCI-3242, CCI-3243, CCI-3244, CCI-3245, CCI-3246

The \href{http://isocpp.github.io/CppCoreGuidelines/CppCoreGuidelines#S-functions}{C++ Core Guidelines} are used as the standard for development. The \href{https://google.github.io/styleguide/cppguide.html}{Google C++ Style Guide}, \href{https://resources.sei.cmu.edu/downloads/secure-coding/assets/sei-cert-cpp-coding-standard-2016-v01.pdf}{CERT Coding Standards}, and \href{https://wiki.qt.io/Qt_Coding_Style}{Qt Coding Style} are used as secondary references.

Since there is a single user role for the application, the design document is defined as the User's Guide.

Security requirements are to perform the testing listed in Section~\ref{sec:sa-11}. Changes to these requirements are stored in the GitHub changelog for each file. Reviews of this process occur for every major release.

\subsubsection{SA-15(4) -- Threat Modeling / Vulnerability Analysis}

\paragraph{Applicable CCIs:} CCI-3256, CCI-3257, CCI-3258, CCI-3259, CCI-3260, CCI-3261, CCI-3262, CCI-3263, CCI-3264, CCI-3265, CCI-3266, CCI-3267, CCI-3268, CCI-3269, CCI-3270, CCI-3271

Vulnerability analysis includes the scans listed in Section~\ref{sec:sa-11}.

The following sections are required by STIG rule SV-85011r1\_rule:
\begin{enumerate}
	\item \textbf{Identified Threats:} none
	\item \textbf{Potential Vulnerabilities:} XML Bomb (Lots of Laughs attack, \href{https://github.com/squinky86/STIGQter/issues/26}{issue \#26})
	\item \textbf{Counter Measures Taken:} none; The threat of malicious CKL files is well-documented.
	\item \textbf{Potential Mitigations:} Train the user not to open CKL files from unknown sources.
	\item \textbf{Mitigations Selected based on Risk Analysis:} A risk acceptance is recommended.
\end{enumerate}

\subsubsection{SA-15(5) -- Attack Surface Reduction}

\paragraph{Applicable CCIs:} CCI-3272

Error handling vulnerabilities are included in the scanning process of Section~\ref{sec:sa-11}.

\subsubsection{SA-15(10) -- Incident Response Plan}

\paragraph{Applicable CCIs:} CCI-3289

All incidents should be logged as issues in the GitHub project's issue tracker. These incidents will be triaged on a best-effort basis.

\subsection{SA-22 -- Unsupported System Components}

\paragraph{Applicable CCIs:} CCI-3374, CCI-3376

When a dependency is updated for a security reason, the build revision of STIGQter is automatically updated as well to enforce updates. This includes vulnerabilities in any of the dependent components. For example, if STIGQter 1.0.0 is released and the next day a Qt vulnerability is released that affects STIGQter, STIGQter documentation will be updated and a new release of STIGQter 1.0.1 will occur. New releases are made available on https://www.stigqter.com/.

When no longer supported, the STIGQter website will be updated to reflect its status.

\paragraph{Non-Applicable CCI(s):} CCI-3375

There is no justification to use outdated components. Only supported products and dependencies are accepted.

\clearpage
\printbibliography

\end{document}
